\documentclass[main.tex]{subfiles}

\begin{document}

\section{Bilan du projet}

\subsection{Apports individuels et collectifs}

\quad La réalisation de ce projet a représenté une opportunité concrète de mettre en œuvre les compétences acquises durant notre formation en BTS CIEL, autant sur le plan individuel que sur le plan collectif.

Chaque membre de l’équipe a pu approfondir ses connaissances techniques en lien avec son domaine d’intervention.

\begin{itemize}

\item L’élève Antoine Canin a renforcé ses compétences en configuration de réseau (IPFire), en installation de base de données, ainsi qu’en intégration de lecteur RFID.
\item L'élève Alexandre Picard s’est spécialisé dans le développement logiciel sous C\#, dans la gestion de base de données, dans le chiffrement de données, ainsi que dans le design applicatif.
\item L’élève William Pinalie a développé ses compétences en électronique embarquée, plus particulièrement dans le pilotage de servomoteurs et de diodes RGB avec une Raspberry Pi grâce au langage C++.

\end{itemize}

\subsection{Conclusion générale}

\quad Le projet a été mené à bien dans les délais impartis avec un prototype fonctionnel et conforme au cahier des charges initial. Il répond à une problématique réelle de l’entreprise Sylvamo en apportant une solution innovante et automatisée pour la gestion des documents d’identités des chauffeurs routiers. 

Sur le plan technique, l’intégration entre le matériel (Lecteur RFID, Servomoteurs, Raspberry PI, NUC) et du logiciel (Application Windows, Base de données MariaDB, Pare-feu IPFire) a été réalisée avec succès. Le système est évolutif , sécurisé et conçu pour une utilisation en environnement industriel.

Ce projet a donc représenté une expérience professionnelle riche et formatrice, nous préparant aux exigences du monde de l’entreprise. Il a également renforcé notre capacité à travailler en équipe autour d’un objectif commun tout en valorisant les compétences individuelles de chacun.



\newpage

\end{document}